\Needspace{13\baselineskip}
\chapterrule
\section{数据、训练系统与部署：资源怎样分配？}
\label{intro:23b-deployment-and-ecosystem}



本章把生产链分成三段：先把资料整理成训练样本，再分配训练计算与通信，
最后安排推理时的权重、缓存、延迟与成本。
第~\ref*{intro:23a-from-toy-to-gpt}~章讨论了模型规模，这里先接上它的数据来源；
第~\ref*{intro:23d-agent-tools}--\ref*{intro:23d-agent-protocols}~章搭好的助手，则在最后一段落实服务预算。

\subsection{训练数据工程：来源、清洗、去重与配比}
\label{sec:data-engineering}
\phantomsection\label{concept:pretraining-data}
\phantomsection\label{concept:evaluation-contamination}
预训练先把网页、书籍、代码与问答整理成\textbf{来源可追溯的 Token 样本}。

\begin{enumerate}[itemsep=0pt,parsep=0pt,topsep=2pt]
\item \textbf{来源与授权：}记录许可和使用范围，剔除敏感信息与禁止使用的数据；
\item \textbf{清洗：}去掉破损编码、导航和广告模板，保留正文与必要上下文；
\item \textbf{去重：}哈希识别完全重复，分片相似度发现近重复候选；
\item \textbf{质量与配比：}平衡语言和领域，测量各来源贡献，避免模板挤占信号；
\item \textbf{切分与污染：}按组划分训练／验证／测试，核对测试题和答案是否泄漏。
\end{enumerate}

图中四条文本先做全半角规范化，并去掉空格和标点。A、B 的正文与完整 SHA-256 相同，保留一份；图中显示哈希前八位。

\par\Needspace{75mm}
\tagstructbegin{tag=Figure,alttext={四条小文本经规范化后：A 与 B 都是春天花开了，SHA256 前八位同为 c0ac0b65，合并精确重复；C 是春天的花开了，哈希不同但与 A 的二字片段 Jaccard 相似度为二分之一，进入近重复复核；D 是春天花谢了，相似度三分之一，且事实相反，应保留。}}
\begin{center}
\begin{tikzpicture}[>=Latex,x=1mm,y=1mm,
 n/.style={draw=linegray,rounded corners=2pt,fill=paper,align=left,
 font=\small,minimum height=12mm,text width=58mm}]
\node[n] (a) at (0,0) {A：春天花开了\\B：春天 花 开了。};
\node[n,fill=paper] (ab) at (78,0) {规范化后完全相同\\哈希 \texttt{c0ac0b65}，保留一份};
\node[n] (c) at (0,-16) {C：春天的花开了\\哈希 \texttt{db2d9197}};
\node[n,fill=linegray!25] (near) at (78,-16) {$J(A,C)=3/6=0.5$\\近重复候选 $\to$ 内容复核};
\node[n] (d) at (0,-32) {D：春天花谢了\\哈希 \texttt{b3dc2f3b}};
\node[n,fill=linegray!15] (keep) at (78,-32) {$J(A,D)=2/6\approx0.333$\\含义相反，保留不同事实};
\draw[->,thick,flowcolor] (a)--(ab);
\draw[->,thick,flowcolor] (c)--(near);
\draw[->,thick,flowcolor] (d)--(keep);
\end{tikzpicture}
\captionof{figure}{先识别完全重复，再复核近重复；相似文本中的事实差异值得保留}
\label{fig:corpus-dedup}
\end{center}
\tagstructend


C 与 A 的哈希不同，再把正文切成相邻二字片段的集合：
\[
\begin{aligned}
S_A&=\{\text{春天、天花、花开、开了}\},\\
S_C&=\{\text{春天、天的、的花、花开、开了}\}.
\end{aligned}
\]
两者交集有三个片段、并集有六个，Jaccard 相似度为
$J(A,C)=|S_A\cap S_C|/|S_A\cup S_C|=3/6=0.5$。
本例复核 $J\ge0.5$ 的候选，确认同源后合并；D 的分数为 $2/6$，且“谢”与“开”表达相反事实，应保留。

运行 \texttt{python3 scripts/pretraining\_data\_demo.py} 可复算全部哈希与相似度。
同源文本先分组，再按组切分训练／验证／测试集，避免同一内容跨集合泄漏。
若测试题、答案或它们的近重复版本进入了训练集、微调集、检索库，或在反复调参中被看过，
评测分数就会混入对见过内容的记忆，这称为\textbf{评测污染}。
上面的哈希与片段相似度同样可用来比对测试集；在可行时再用时间隔离的保留集与来源追踪，并对新题复测。

\textbf{应用条件：}本例删除标点、采用二字片段和 $0.5$ 阈值。
真实语料需保留有意义的符号，按领域选择分片与阈值；金额、日期或否定词的变化须内容复核。

\subsection{训练系统：参数、梯度、状态与通信}
\subsubsection*{为什么大模型训练要拆到多张卡？}

权重只是训练显存的一部分。还要保存梯度、优化器状态和
反向传播需要的激活。以十亿参数为教学估算：
FP32 表示每个数用 $32$ bit，也就是 $4$ 字节保存。
仅主权重、梯度与 Adam 为更新参数保存的两份历史状态，共四份十亿个数，
理论上就需要 $10^9\times4\times4=16\times10^9$ 字节，即约 $16$ GB，
激活、临时缓冲和额外权重副本需另计。

\subsubsection*{激活显存与梯度检查点：用重算换空间}
参数量相同，序列或批次变长，训练显存也可能显著增加。
反向传播需要前向的中间结果，如每层输入、Q/K/V 与 FFN 激活；
这些\textbf{激活}随当前批次变化，不属于模型参数。
设批次大小 $B=2$、长度 $T=2048$、宽度 $d=1024$，
一张 BF16 隐藏表占 $BTd\times2=8{,}388{,}608$ 字节，即 $8$ MiB。
若 $24$ 层每层只保留这样一张表，就已是 $192$ MiB，实际还会保存更多中间量。
朴素实现若以每数 $2$ 字节另存每层 $H=16$ 个头的完整注意力表，
单层就需 $BHT^2\times2=256$ MiB；长度翻倍时，这一项变为四倍。
是否物化这张表、采用什么内核，会改变实际占用。

\phantomsection\label{concept:activation-checkpoint}%
\textbf{梯度检查点（激活重计算）}只保存选定的层边界输入，
反向到达某段时再运行该段前向，重建需要的中间结果。
例如上述 $24$ 层每四层一段，六份段入口表共 $6\times8=48$ MiB，
在反向时逐段重算；还需当前段的中间激活与工作空间，所以 $48$ MiB 只是入口表的分项。
重算要复现原前向的随机状态等条件，才能得到对应梯度。
它减少激活保存量，参数、梯度和优化器状态仍按原方案存储；
方法来源见\hyperref[source:activation-checkpoint]{激活重计算论文}。

\subsubsection*{优化设置也影响单位计算能学到多少}
主线用固定学习率的梯度下降来显示每项更新。
大规模训练还需记录\textbf{学习率调度}：Warmup 在开头逐步提高学习率，
之后按选定曲线衰减；例如先从 $10^{-5}$ 逐步升到 $10^{-3}$，再逐渐降到 $10^{-4}$。
调度需要与批次、总步数和损失曲线一起检查。
Adam 用梯度的一阶与二阶历史调整更新，AdamW 则把权重衰减与自适应梯度项分开处理；
完整手算见\hyperref[app:optimizer-details]{附录 A 的优化器选读}，
解耦依据见\hyperref[source:adamw]{AdamW 原论文}。

\subsubsection*{再按瓶颈选择并行方式}
多张设备共同完成训练，叫作\textbf{分布式训练}。怎样分，要先看遇到哪个瓶颈。

\textbf{放得下，但样本算得慢：先分数据。}
两张卡各自保存完整模型，各处理一半样本，再汇总梯度、更新为相同的新权重。
等大批次的梯度取平均，就对应两批合成的 Mini-batch 梯度。
这叫\textbf{数据并行}：分摊了样本计算，朴素方案每卡仍需上述约 $16$ GB 的训练状态，激活另计。

\textbf{放不下：看能把哪部分拆开。}
若主要是权重、梯度和优化器状态占空间，可以把它们分片，计算需要时再通信取用，
这叫\textbf{状态分片}。若四份状态全部均分到两卡，理想的持久存储约为每卡 $8$ GB；
加上激活、临时聚合和通信缓冲，每卡的总显存需求还会高于 $8$ GB。
也可以把四层网络的前两层放卡 1、后两层放卡 2，传递中间激活，
这叫\textbf{流水线并行}；多份微批次依次进入，才能让不同阶段同时工作。
若连一层的大矩阵都需要拆，则可把一个 $4\times4$ 权重按输出列分成两张 $4\times2$ 表，
两卡分别算一半输出，再拼接，这叫\textbf{张量并行}。

\par\noindent\begin{minipage}{\linewidth}
\tagstructbegin{tag=Figure,alttext={同用两张卡的四种分工。数据并行各有完整模型，分别处理样本并同步梯度；状态分片各保存半份状态并按需通信；流水线卡一保存第一二层、卡二保存第三四层并传激活；张量并行各保存同一矩阵一半输出列并拼接结果。}}
\begin{center}
\begin{tikzpicture}[x=1mm,y=1mm,>=Latex,
 device/.style={draw=linegray,rounded corners=2pt,fill=paper,align=center,font=\footnotesize,text width=39mm,minimum height=12mm}]
\node[font=\small\bfseries] at (45,10) {卡 1};
\node[font=\small\bfseries] at (111,10) {卡 2};
\node[font=\footnotesize,anchor=east] at (19,0) {数据并行};
\node[device] (d1) at (45,0) {完整模型\\样本前半批};
\node[device] (d2) at (111,0) {完整模型\\样本后半批};
\draw[<->,thick] (d1)--node[above,font=\scriptsize]{梯度}(d2);
\node[font=\footnotesize,anchor=east] at (19,-19) {状态分片};
\node[device] (s1) at (45,-19) {状态分片 1};
\node[device] (s2) at (111,-19) {状态分片 2};
\draw[<->,thick] (s1)--node[above,font=\scriptsize]{按需取用}(s2);
\node[font=\footnotesize,anchor=east] at (19,-38) {流水线并行};
\node[device] (p1) at (45,-38) {第 1、2 层};
\node[device] (p2) at (111,-38) {第 3、4 层};
\draw[->,thick] (p1)--node[above,font=\scriptsize]{激活}(p2);
\node[font=\footnotesize,anchor=east] at (19,-57) {张量并行};
\node[device] (t1) at (45,-57) {同一层权重第 1、2 列};
\node[device] (t2) at (111,-57) {同一层权重第 3、4 列};
\draw[<->,thick] (t1)--node[above,font=\scriptsize]{拼接输出}(t2);
\end{tikzpicture}
\captionof{figure}{先问拆什么：数据、训练状态、完整层，还是一层中的矩阵}
\label{fig:distributed-two-devices}
\end{center}
\tagstructend
\end{minipage}\par


\begin{center}\small
\tagpdfsetup{table/header-rows={1}}
\begin{tabularx}{\linewidth}{@{}p{31mm}p{35mm}>{\raggedright\arraybackslash}X@{}}
\toprule
做法 & 主要分摊对象 & 主要代价\\\midrule
数据并行 & 样本计算 & 需同步梯度，完整模型仍需放得下\\
状态分片 & 优化器状态、梯度、权重 & 按分片范围通信取用与归约，如 ZeRO/FSDP\\
流水线并行 & 完整层 & 传递激活，等待会形成空泡\\
张量并行 & 一层中的矩阵与计算 & 层内需要频繁通信\\
\bottomrule
\end{tabularx}
\end{center}
先找到瓶颈，再选择分摊对象；实际训练可以组合使用这些办法。
对应实例见附录 F：状态分片有\hyperref[source:zero]{ZeRO}与\hyperref[source:pytorch-fsdp]{PyTorch FSDP}，
流水线并行见\hyperref[source:gpipe]{GPipe}，层内张量并行见\hyperref[source:megatron-lm]{Megatron-LM}。

\phantomsection\label{concept:mixed-precision}
\textbf{一个数占多少空间，也可以调整。}
若某项数据从每数 $32$ bit 改为 $16$ bit，该项存储就从每数 $4$ 字节减为 $2$ 字节。
\textbf{混合精度训练（Mixed Precision）}对适合的矩阵计算使用如 BF16 的较短浮点表示，
对某些归约、统计或状态保留较高精度，以兼顾计算成本和数值稳定性。
浮点数用“有效数字乘以二的某次幂”表示数值：指数位关系到能表示多大或多小，尾数位关系到数值有多精细。
BF16 具有与 FP32 相同宽度的指数位，因此动态范围较大，
但尾数精度更低；采用哪种精度取决于硬件及数值稳定性。
混合精度训练的一种具体实现还维护 FP32 主权重并使用损失缩放，见\hyperref[source:mixed-precision]{附录 F 的论文来源}。

\textbf{还有一种慢，来自计算与显存之间反复搬运数据。}
标准 Attention 若把整张分数和权重矩阵写入显存再读出，会产生大量中间数据读写。
\textbf{FlashAttention} 通过分块和逐块累积的 Softmax 计算，减少这类读写与中间存储，
提高执行效率；它在适用条件下\textbf{保持相同的注意力数学目标}。
其分块算法与显存读写分析见\hyperref[source:flash-attention]{附录 F 的论文来源}。
工程选型要同时测量吞吐、显存、通信、数值误差与最终模型质量。

\subsection{推理部署：权重、缓存、精度与服务}
\subsubsection*{训练代码、模型权重、推理程序是三件事}

让一个大语言模型工作，至少需要区分：
\begin{enumerate}
\item \textbf{模型结构与代码}：决定矩阵和非线性如何排列、计算。
\item \textbf{模型权重}：训练得到的大量参数数值。
\item \textbf{推理运行时}：加载权重、完成 Tokenization、执行网络计算并生成输出的程序。
\end{enumerate}
有结构和代码但没有训练好的权重，模型还无法提供成熟的语言能力；
拿到权重后，也需要合适的推理程序才能运行。

\textbf{模型}是有明确架构、参数和版本的计算系统；\textbf{产品}是用户接触的应用、界面或 API，
往往还组合检索、工具、审核与计费，所以比较时分别记录产品功能与所用权重版本。

\textbf{开放权重}意味着外部能够获取模型参数；开源程度则应分别核对训练代码、数据、再分发条款与商业许可。

\textbf{私有化部署}一般是指在自有设备或受控制的基础设施上运行模型。它有助于控制数据流向，但并不自动带来安全保障：仍需考虑访问权限、日志、模型许可和升级流程。

\subsubsection*{模型为什么需要显存？}

权重由许多数字组成。假设模型有 $10^9$（十亿）个参数，如果每个参数用 2 字节存储，仅权重就占约
\[
10^9\times2\ \text{bytes}\approx2\ \text{GB}.
\]
这是十进制粗略估计，还不包含 \conceptlink{kv-cache}{KV Cache}、中间激活、运行时缓冲、框架开销等。

\textbf{把权重和并发缓存放进同一笔预算。}
假设一项服务使用 $8\times10^9$ 个参数、FP16 权重，模型有 $32$ 层，每层 $8$ 个 K/V 头、每头 $128$ 维；
每个请求最多保留 $8192$ 个位置，K 和 V 均用 $2$ 字节存储。
权重约占
\[
8\times10^9\times2\ \text{B}=16\times10^9\ \text{B}\approx14.9\ \text{GiB}.
\]
单个请求的 KV Cache 本体为
\[
\begin{aligned}
2_{\text{K,V}}\times32_{\text{层}}\times8192_{\text{位置}}
\times8_{\text{K/V 头}}\times128_{\text{头维}}\times2_{\text{B/数}}
&=1{,}073{,}741{,}824\ \text{B}\\
&=1\ \text{GiB}.
\end{aligned}
\]
若同时服务 $4$ 个都达到该长度的请求，仅权重与缓存就约需 $14.9+4=18.9$ GiB；运行时缓冲、激活和显存分配开销还要另留空间。
这是容量下限估算，不等于设备的实际峰值显存。

服务预算要区分阶段测量：Prefill 记录输入处理耗时或吞吐，Decode 记录生成阶段的每 Token 延迟或 Token 吞吐；首 Token 延迟还包含排队、Prefill 和首步生成。固定硬件、精度、输入／输出长度和并发，实测这些指标与峰值显存；它们不能从参数量或缓存公式直接推出。

\subsubsection*{服务调度：谁先算，缓存放哪里？}
\phantomsection\label{concept:continuous-batching}%
\textbf{连续批处理}在每轮生成后移走已完成的请求，再将新请求加入可用位置。
例如一批 A、B 分别还需生成 $1$、$4$ 个 Token，C 已在等待：
第一轮 A 完成后，就可接纳 C 的 Prefill；C 准备好后与 B 一同 Decode，
无须等 B 全部结束。
这是请求调度的变化；每条请求内部仍按自身前缀自回归生成。
长 Prompt 的 Prefill 还可能占住设备，使正在 Decode 的请求等待，
可将 Prefill 分块并设置每轮 Token 预算，在首 Token 延迟与生成间隔之间取舍，
见\hyperref[source:sarathi-serve]{Sarathi-Serve 的调度方案}。

\phantomsection\label{concept:paged-attention}%
\textbf{PagedAttention}把 KV Cache 分成固定大小的块，
通过块表把请求的逻辑位置映射到可分散存放的物理块，按需分配与回收。
例如每块放 $4$ 个位置，长度为 $9$ 的请求占 $3$ 块，最后一块尚有 $3$ 个空位；
无需从一开始预留整条最大长度的连续缓存。
它改善缓存的分配和碎片问题，每个已缓存位置的 K/V 数值仍需空间。
连续批处理与分页缓存的配合见\hyperref[source:paged-attention]{vLLM/PagedAttention 论文}。

\phantomsection\label{concept:speculative-decoding}%
\textbf{投机解码}先由较小的草稿模型连续提出若干 Token，
再由目标模型一次并行计算这些候选位置的概率，并按接受规则保留一个前缀。
若在中途拒绝，就丢弃后面的草稿，并从校正后的分布补出一个 Token。
采用\hyperref[source:speculative-decoding]{原论文的精确接受与校正规则}时，
可以保持目标采样分布；直接照收草稿或只比较最高分，得不到这个结论。
收益取决于接受率、草稿开销和批处理负载，还需额外存放草稿模型。
这些方案分别改变调度、缓存管理和解码计算，评测时应同时报告吞吐、
首 Token 延迟及生成间隔的分位数，检查高负载下的等待是否可接受。

\textbf{GPU} 擅长对大量数字并行计算；\textbf{TPU} 是面向机器学习工作负载设计的专用加速器；\textbf{NPU} 泛指神经网络计算加速单元。这些硬件为模型提供执行矩阵运算和数据搬运的计算基础设施。

\begin{warningbox}
训练时还常需保存 \conceptlink{gradient}{梯度}、优化器状态等额外数据，这些都要在权重存储之外另计显存。即使模型权重能放入显存，运行长上下文或高并发时仍可能发生内存不足。

硬件名称说明设备类别；算子与精度支持、带宽和延迟须按具体型号、运行时与负载核对。
\end{warningbox}


\subsubsection*{量化：为什么更少的位数可以降低成本？}
\phantomsection\label{concept:quant}

原来每个权重使用较高数值精度保存；\textbf{量化}用更少位数表示近似的权重或激活。例如从每个参数 16 bit 变为 8 bit，权重存储理论上可下降到原来的一半左右。

\begin{example}
\noindent\emph{把四个权重实际量化到 8 bit}\par

取浮点权重 $(-1,-0.2,0.2,1)$，
使用包含零点的\textbf{对称有符号 8-bit 量化}，
整数范围 $[-127,127]$，比例尺 $s=1/127$。
将每个数除以 $s$、四舍五入并截取合法范围：
\[
(-1,-0.2,0.2,1)
\longrightarrow(-127,-25,25,127).
\]
反量化（整数乘回 $s$）得到约
$(-1,-0.197,0.197,1)$。
存储从每数 16 bit 变为每数 8 bit 时，
仅这四个权重从 8 字节降为 4 字节，不包括比例尺和分组元数据。
\end{example}

代价是近似误差：相近权重会落入同一个量化格。模型效果是否受影响，取决于量化方案、模型分布和任务，并不是位数越少越好。
整数算术推理的一种实现路线见\hyperref[source:integer-quantization]{附录 F 的量化论文}。

\subsubsection*{\texorpdfstring{{\symboliconfont ★}}{★} 进一步阅读：蒸馏与 LoRA}

如果希望训练一个较小模型承担任务，可以让较强模型提供学习信号。
\phantomsection\label{concept:distillation}%
这种利用教师模型指导学生模型的办法叫作\textbf{蒸馏}。
教师除了给出答案，还能给出各候选的概率：

\begin{example}
\noindent\emph{硬标签与软标签怎样不同？}\par

教师模型给两个类别的概率为 $(0.8,0.2)$。
硬标签只告诉学生“类别 1 是正确答案”，可表示成 $(1,0)$；
\textbf{软标签}则告诉学生：类别 2 仍有 $0.2$ 的概率。
如果学生预测 $(0.6,0.4)$，硬标签交叉熵为
$-\ln0.6\approx0.511$，
而对教师软标签的交叉熵为
\[
-\bigl(0.8\ln0.6+0.2\ln0.4\bigr)\approx0.592.
\]
后者同时考虑两个类别的匹配程度。蒸馏时还可以调整温度，并把软标签损失与原任务损失结合。
\end{example}

\begin{example}
\noindent\emph{LoRA 的参数量真正省在哪里？}\par

另一种处境是：已有模型能运行，但为了适应任务而更新全部大矩阵，训练开销太大。
能否保留原权重，只学习一份较小的修正？
\phantomsection\label{concept:lora}%
\textbf{LoRA（低秩适配）}冻结原矩阵 $W$，用两个窄矩阵的乘积表示修正 $\Delta W$，
使新的计算成为 $x(W+\Delta W)=xW+x\Delta W$；训练只更新修正分支的参数。

先用独立的二维例看输出怎样改变，暂把缩放系数取为 $1$：
\[
x=(1,2),\quad W=I_2,\quad
U_{\rm lr}=\begin{bmatrix}1\\0\end{bmatrix},\quad V_{\rm lr}=(0,1),\quad
\Delta W=U_{\rm lr}V_{\rm lr}=\begin{bmatrix}0&1\\0&0\end{bmatrix}.
\]
原来 $xW=(1,2)$，修正分支给出 $x\Delta W=(0,1)$，相加后为 $(1,3)$。
原矩阵虽冻结，预测仍能随修正分支改变。再来看大矩阵能节省多少可学习参数。

对一个 $1000\times1000$ 的方阵，完整训练它需要
$1000^2=1{,}000{,}000$ 个可学习权重。
LoRA 若冻结原矩阵，用中间宽度 $r=8$ 的两个窄矩阵表示增量
$\Delta W=U_{\rm lr}V_{\rm lr}$，$U_{\rm lr}:1000\times8$，
$V_{\rm lr}:8\times1000$，则需要
\[
1000\times8+8\times1000
=\boxed{16{,}000}
\]
个可学习参数，约为完整权重的 $1.6\%$。
$r$ 限制修正能独立组合的方向数量，是增量矩阵的\textbf{秩上限}；这里称为低秩，是相对于原矩阵的宽度 $1000$。
这也限制了修正的自由度，效果要随任务与 $r$ 的选择验证。
这指的是该矩阵新训练的增量参数数目，
不是整套模型参数或部署占用必然缩小到 $1.6\%$。
冻结与增量相加的机制见\hyperref[source:lora]{LoRA 原论文}。
\end{example}

\phantomsection\label{concept:pruning}%
若希望直接减少原模型中的权重、连接或结构，还可以评估删去哪些部分，这叫作\textbf{剪枝}。
将几种方法按它们改动的对象归纳如下：
\begin{center}
\small
\renewcommand{\arraystretch}{1.23}
\tagpdfsetup{table/header-rows={1}}
\begin{tabularx}{\linewidth}{@{}p{38mm}X@{}}
\toprule
技术 & 核心目的\\
\midrule
量化（Quantization） & 用低精度数值减少模型存储或计算成本\\
蒸馏（Distillation） & 利用较大或较强模型的输出作为训练信号，帮助学生模型学习\\
剪枝（Pruning） & 去除或稀疏化较不重要的权重、连接或结构\\
LoRA & 在微调阶段用低秩增量参数，减少需要训练的参数量\\
\bottomrule
\end{tabularx}
\end{center}

% Keep this closing block with the self-check so the part-end page is not a sparse tail.
\Needspace{15\baselineskip}
\subsubsection*{从权重到服务：工具链各管哪一段？}

常见工作链是：模型权重 $\to$ 推理引擎 $\to$ 接口服务 $\to$ 应用与工具。
\textbf{框架}执行训练与推理，\textbf{镜像}打包环境，\textbf{推理引擎}安排批处理与缓存；
\textbf{接口和协议}约定 \conceptlink{agent}{Agent} 怎样请求模型、通过 \conceptlink{mcp}{MCP} 接入工具。
本书的推理过程固定模型参数。

\begin{center}\small
\tagpdfsetup{table/header-rows={1}}
\begin{tabularx}{\linewidth}{@{}p{35mm}X@{}}\toprule
按需求查阅 & 所处环节\\\midrule
\href{https://docs.nvidia.com/cuda/cuda-programming-guide/index.html}{CUDA}、\href{https://docs.pytorch.org/docs/stable/index.html}{PyTorch} & 硬件计算平台与张量框架，执行训练及矩阵运算\\
\href{https://huggingface.co/docs/hub/index}{Hugging Face Hub} & 保存模型权重与说明\\
\href{https://docs.ollama.com/}{Ollama}、\href{https://docs.vllm.ai/}{vLLM} & 本地运行或服务型推理工具；加载权重，管理执行、批处理或缓存\\\bottomrule
\end{tabularx}
\end{center}
比较推理工具时固定权重、硬件、输入长度、并发数与量化方案。
软件版本、设备兼容与许可证以项目文档为准。

\chapterchecktasks{23b-deployment-and-ecosystem}{理解与复算}{
\item \textbf{区分：}训练代码、模型权重、推理引擎和对外 API 在一个产品里承担什么不同职责？\chapterinlinehint{按训练、存储、执行和服务四个环节追踪职责。}
\item \textbf{估算：}$10^9$ 个参数各占 $2$ 字节，不计运行开销仅权重需要多少十进制 GB？若另有 $0.5$ GB 缓存，合计至少多少？\chapterinlinehint{先算参数个数乘单个字节数，再统一 GB 单位。}
\item \textbf{找错：}有人认为把参数量化到更少位数，模型质量、吞吐与延迟都会按字节比例等幅改善。请指出两处需要实测的地方。\chapterinlinehint{分别设计质量对照与同负载的速度测试。}
}
